L'architettura sviluppata risponde alle domande poste nel progetto: 
\begin{itemize}
    \item Implementa un classificatore valutato con criteri severi e vicini alle condizioni real-time;
    \item Propone un miglioramento per le classiche pipeline attraverso il meccanismo di selezione delle finestre;
    \item Per la costruzione di un predittore sperimenta una metodologia ancora inutilizzata, per questo preciso dataset, basata sulle onde mu.
\end{itemize}

\section{Conclusioni}
Sul piano delle prestazioni, il sistema ottiene in media il 73,61\% nella distinzione fra riposo e attivazione, il 69,98\% fra pugni e piedi e il 62,55\% fra mano sinistra e destra, con quote di finestre scartate comprese fra il 19 e il 26\%; la catena completa a cinque classi raggiunge il 36,46\%, a fronte di un livello del caso del 20\%. Si tratta di valori da leggere alla luce delle condizioni deliberatamente severe in cui sono stati ottenuti: soggetti non addestrati, nessun riscontro durante la registrazione, due sole run di addestramento, validazione su run mai viste e decisione su finestre scorrevoli anziché su trial ritagliati.

La media, tuttavia, descrive solo in parte il comportamento del sistema. Il 10\% migliore della popolazione supera in media il 94\% nel compito pugni/piedi scartando appena il 4\% delle finestre, e si attesta intorno all'89\% anche nel compito sinistra/destra, il più arduo; nella catena a cinque classi il soggetto migliore riconosce correttamente tre finestre su quattro. Per una parte dei soggetti la pipeline raggiunge dunque accuratezze vicine a quelle di un sistema utilizzabile nella pratica, pur in assenza di qualunque addestramento; all'estremo opposto, meno della metà dei soggetti supera la soglia di significatività nel compito sinistra/destra.

Considerati nel loro insieme, tutti gli elementi trattati indicano che, \textbf{entro i limiti del dataset e della pipeline considerati, la variabilità inter-soggetto emerge come uno dei principali fattori associati alla prestazione}: le scelte di progetto esaminate spostano la media di pochi punti percentuali, mentre la distanza fra i soggetti supera i sessanta punti e si conserva attraverso compiti e registrazioni distinti.

Occorre tuttavia precisare la portata di questa affermazione, poiché i risultati non consentono di separare in modo netto il contributo della capacità fisiologica della persona da quello degli altri fattori in gioco. In primo luogo, l'architettura e i suoi parametri (sotto-bande del banco di filtri, durata delle finestre, canali, soglia probabilistica) sono identici per tutti i soggetti, e il solo iperparametro calibrato individualmente è il coefficiente di contrazione: una configurazione poco adatta a un particolare soggetto si manifesta quindi come una bassa prestazione di quel soggetto. In secondo luogo, tutte le registrazioni di un soggetto provengono da un'unica sessione, per cui la coerenza fra compiti può riflettere anche fattori condivisi dalla sessione, come la qualità del contatto degli elettrodi, il livello di attenzione o la strategia di immaginazione adottata. In terzo luogo, come osservato nel capitolo sul dataset, un soggetto incapace di generare un segnale utilizzabile e un soggetto semplicemente non addestrato sono in questo contesto indistinguibili, e la scarsità dei dati di addestramento, appena due run per modello, penalizza in misura maggiore proprio i soggetti con un segnale meno marcato. Il predittore, che fra le grandezze esaminate è l'unica misurata indipendentemente dal classificatore, spiega del resto soltanto una frazione limitata della variabilità osservata.

Ciò che i risultati permettono di affermare è dunque che una quota rilevante della prestazione è legata a caratteristiche del soggetto stabili almeno all'interno della sessione, e che tali caratteristiche sono in parte riconoscibili già dall'attività a riposo. Stabilire quanto di essa costituisca una proprietà fisiologica della persona e quanto dipenda invece dall'addestramento, dalla sessione o dalla configurazione del classificatore richiederebbe registrazioni ripetute nel tempo e un'ottimizzazione specifica per soggetto, che esulano dal perimetro di questo lavoro.


\section{Limiti e sviluppi futuri}
Ogni sperimentazione effettuata rappresenta una solida base, con le dovute limitazioni, da cui poter derivare modelli con maggiore complessità e precisione. 

L'approccio del classificatore basato sulla geometria riemanniana, grazie all'applicazione del riallineamento, si presta ad adattamenti per valutazioni cross-subject o applicazioni di transfer learning. La presenza di criteri valutativi rigidi come la LORO e la k-fold cross validation, affiancati alla metodologia delle sliding windows (meno generosa del tipico epoching), fornisce un modello rigido ma con elevata interpretabilità e affidabilità.

La selezione delle finestre attualmente non rappresenta un approccio fortemente migliorativo: l'aumento medio dell'accuratezza non compensa adeguatamente il numero di classificazioni scartate. Nonostante ciò è importante sottolineare, come già osservato nel capitolo precedente, la differenza che intercorre tra i soggetti a più alte prestazioni e quelli affetti da illiteracy o che necessitano addestramento. Per la prima categoria, a fronte di un modesto scarto è possibile apprezzare un miglioramento delle performance. Il prossimo step potrebbe dunque prevedere lo sviluppo di un algoritmo specializzato che, per ogni soggetto, individui il rapporto ottimale tra scarti e miglioramento. Completato questo passaggio sarà possibile sperimentare l'approccio attraverso un sistema real-time: valutando la percezione e i feedback degli utenti e analizzando parametri quali il numero di finestre scartate per ogni compito sarà possibile ottenere, dallo stesso sistema di partenza, un comportamento più generoso per l'utente finale.

Infine, il predittore di illiteracy fornisce già allo stato attuale una correlazione significativa, seppur moderata, tra i ritmi a riposo e le performance relative, ma si limita ad osservare un singolo minuto di run per soggetti non addestrati. Studi più approfonditi potrebbero analizzare l'andamento del picco mu a seguito di una serie di sessioni di addestramento o affiancare altri parametri valutativi e mescolare le valutazioni. Inoltre, allo stato attuale, essendo il coefficiente di Spearman non particolarmente elevato, non è stato addestrato nessun modello in grado di distinguere soggetti a basse performance da soggetti ad alte performance.